\newcommand{\X}{\mathcal{X}}
\newcommand{\I}{\mathcal{I}}
\newcommand{\J}{\mathcal{J}}
%\newcommand{\E}{\mathbb{E}}
\renewcommand{\P}{\mathcal{P}}
%\newcommand{\Yobs}{Y}
\newcommand{\Yobst}{Y^{obs,t}}
\newcommand{\Yobsc}{Y^{obs,c}}
\newcommand{\var}{\mbox{Var}}
\newcommand{\sd}{\mbox{Sd}}
\newcommand{\cov}{\mbox{Cov}}
%\newcommand{\R}{\mathbb{R}}
\newcommand{\Dt}{\tilde D}
\newcommand{\taun}{\hat{\tau}^{mn}}
\newcommand{\theoremfont}{\itshape}
\renewcommand{\descriptionlabel}[1]{\hspace{\labelsep}\textbf{\descriptionlabelfont #1}}
\newcommand{\D}{\mathcal{D}}
\newcommand{\trace}{\mbox{trace}}
\newcommand{\IMSE}{\mbox{IMSE}}
\newcommand{\EMSE}{\mbox{EMSE}}
\newcommand{\MSE}{\mbox{MSE}}
\newcommand{\BIAS}{\mbox{BIAS}}
%\newcommand{\Pcal}{\mathcal{P}}
\newcommand{\mae}{\gamma_{\scriptsize\mbox{max}}}
\newcommand{\mie}{\gamma_{\scriptsize\mbox{min}}}
\newcommand{\mis}{s_{\scriptsize\mbox{min}}}
\newcommand{\ED}{\mathbb{E}_{\mathcal{D}}}
\newcommand{\Ex}{\mathbb{E}_{x}}
\newcommand{\Pobs}{\mathcal{P}_{\mbox{\tiny observed}}}
%\newcommand{\define}{:=}
\newcommand{\deltat}{\tilde \delta}
%\newcommand{\define}{\overset{\mbox{def}}{=}}
%\newcommand{\emin}{e_{\min}}
%\newcommand{\emax}{e_{\max}}
\newcommand{\Sb}{{\bar S}}
\section{Pseudo Code for CATE Estimators}
\label{app:learners}

In this section, we will present pseudo code for the CATE estimators in this paper. We present code for the meta learning algorithms in Section \ref{app:transfer.algos}.
We denote by $Y^0$ and $Y^1$ the observed outcomes for the control and the treated group. For example, $Y^1_i$ is the observed outcome of the $i$th unit in the treated group. $X^0$ and $X^1$ are the features of the control and treated units, and hence, $X^1_i$  corresponds to the feature vector of the $i$th unit in the treated group. $M_k(Y\sim X)$ is  the notation for a regression estimator, which estimates $x \mapsto \E[Y|X=x]$. It can be any regression/machine learning estimator, but in this paper we only choose it to be a neural network or random forest.

% cateestimate T-learner    	
\begin{algorithm}[H] 
	\begin{algorithmic}[1]
		\Procedure{T--Learner}{$X, \Yobs, W$}
		
		\State $\hat \mu_0 	= M_0(Y^0 \sim X^0)$
		\State $\hat \mu_1	 = M_1(Y^1 \sim X^1)$
		
		\item[]
		\State $\hat \tau(x) = \hat \mu_1(x) - \hat \mu_0(x)$ 
		\EndProcedure
	\end{algorithmic}
	\caption{T-learner}
	
	\label{algo:Tlearner}
		\algcomment{$M_0$ and $M_1$ are here some, possibly different machine learning/regression algorithms.}	
\end{algorithm}

% cateestimate S-learner
\begin{algorithm}[H] 
	\begin{algorithmic}[1]
		\Procedure{S--Learner}{$X, \Yobs, W$}
		
		\State $\hat \mu 	= M(\Yobs \sim (X, W))$
%				\Comment{Estimate $\mu_1$ and $\mu_0$}:
						
		\State $\hat \tau(x) = \hat \mu(x,1) - \hat \mu(x,0)$ 
%				\Comment{Difference is the CATE estimator}:
		\EndProcedure
	\end{algorithmic}
	\caption{S-learner}
	
	\label{algo:Slearner}
	\algcomment{$M(\Yobs \sim (X, W))$ is the notation for estimating $(x,w) \mapsto \E[Y|X=x, W=w]$ while treating $W$ as a 0,1--valued feature.}
\end{algorithm}

\begin{algorithm}[H]
	\begin{algorithmic}[1]	
		\Procedure{X--Learner}{$X, \Yobs, W, g$}
		
		
		\item[]
		\State $\hat \mu_0	= M_1(Y^0 \sim X^0)$ 
				\Comment{Estimate response function}
		\State $\hat \mu_1	 = M_2(Y^1 \sim X^1)$
		
		\item[]
		\State $\Dt^1_{i} =         Y^1_{i}                         -     \hat \mu_0(X^1_{i})$
				\Comment{Compute imputed treatment effects}
		\State $\Dt^0_{i} =     \hat \mu_1(X^0_{i})     -     Y^0_{i}$
		
		\item[]
		\State $\hat \tau_1      = M_3(\Dt^{1}      \sim     X^1)$
				\Comment{Estimate CATE for treated and control}
		\State $\hat \tau_0     = M_4(\Dt^0     \sim     X^0)$
		
		\item[]
		\State $\hat \tau(x) = g(x) \hat \tau_0(x) + (1- g(x)) \hat \tau_1(x)$ 
				\Comment{Average the estimates}
		
		\EndProcedure
	\end{algorithmic}
	\caption[test]{X--learner}

	\label{aglo:Xlearner}
	
	\algcomment{$g(x) \in [0,1]$ is a weighing function which is chosen to minimize the variance of $\hat \tau(x)$. It is sometimes possible to estimate $\cov(\tau_0(x) , \tau_1(x) )$, and compute the best $g$ based on this estimate. However, we have made good experiences by choosing $g$ to be an estimate of the propensity score, but also choosing it to be constant and equal to the ratio of treated units usually leads to a good estimator of the CATE. }	
\end{algorithm}


\begin{algorithm}[t]
  \begin{algorithmic}[1]
  \If{$W_i == 0$}
  \State Update the network $\pi_{\theta_0}$ to predict $Y_i^{obs}$
  \State Update the network $\pi_{\theta_1}$ to predict $Y_i^{obs} + \pi_\tau(X_i)$
  \State Update the network $\pi_\tau$ to predict $\pi_{\theta_1}(X_i) - Y_i^{obs}$
  \EndIf
  \If{$W_i == 1$}
  \State Update the network $\pi_{\theta_0}$ to predict $Y_i^{obs} - \pi_\tau(X_i)$
  \State Update the network $\pi_{\theta_1}$ to predict $Y_i^{obs}$
  \State Update the network $\pi_\tau$ to predict $Y_i^{obs} - \pi_{\theta_0}(X_i)$
  \EndIf
  \end{algorithmic}
  \caption{Y-Learner Pseudo Code}
  \label{algo:Ylearner}
  \label{ylearner}
      \algcomment{This process describes training the Y-Learner for one step given a data point $(Y_i^{obs}, X_i, W_i)$}
\end{algorithm}
